\Aufgabe[15]{Formeln mit Diagrammen darstellen}
{
Diagramme wie im ersten Hefteintrag, die Eingabe, Verarbeitung und Ausgabe darstellen, nennt man Datenflussdiagramm.
\begin{itemize}
    \item Zeichne für eine Wachstumsberechnung und eine Summe aus deiner Tabelle je ein Datenflussdiagramm.
    \item Überlege dabei: Wie stellst du die Daten dar und wieso?\\Zum Beispiel als konkreten Wert, als Zelladresse, als Beschreibung, ... ?
\end{itemize}
\doppelseite{0.48}{0.48}{t}{
    \ifloesung
        \LoesungTikz{
            \node[dfddata, text width=2.2cm, minimum width=3cm, minimum height=1cm] (golf) {Umsatz Golf Q2};
            \node[dfddata, text width=2.5cm, minimum width=3cm, minimum height=1cm, right=0.5cm of golf] (fakt) {Wachstums-faktor};

            \node[oval, text width=2.2cm, minimum width=3.5cm, minimum height=1.2cm, below=0.8cm of $(golf)!0.5!(fakt)$] (stern) {PRODUKT oder *};
            
            \node[dfddata, minimum width=3cm, minimum height=1cm, below=1cm of stern] (gq3) {Umsatz Golf Q3};
            
            \draw[arrow] (golf) -- (stern);
            \draw[arrow] (fakt) -- (stern);
            \draw[arrow] (stern) -- (gq3);
        }
    \else
        \begin{tikzpicture}[
            oval/.style={ellipse, draw, thick, minimum width=3.5cm, minimum height=1.2cm, align=center}
        ]
            \node[dfddata, minimum width=3cm, minimum height=1cm] (golf) {};
            \node[dfddata, minimum width=3cm, minimum height=1cm, right=0.5cm of golf] (fakt) {};
            \node[oval, below=0.8cm of $(golf)!0.5!(fakt)$] (stern) {};
            \node[dfddata, minimum width=3cm, minimum height=1cm, below=1cm of stern] (gq3) {};
        \end{tikzpicture}
    \fi
}{
    \LoesungKaroTikz{
        \node[dfddata, text width=2cm] (gq2) {Umsatz Golf Q2};
        \node[dfddata, text width=2.1cm, right=0.3cm of gq2] (tq2) {Umsatz Tennis Q2};
        \node[dfddata, text width=2.1cm, right=0.3cm of tq2] (sq2) {Umsatz Safari Q2};

        \node[oval, text width=2.2cm, below=0.8cm of tq2] (summe) {SUMME (nicht + ! )};
        
        \node[dfddata, below=1cm of summe] (final) {Gesamtumsatz Q2};
        
        \draw[arrow] (gq2) -- (summe);
        \draw[arrow] (tq2) -- (summe);
        \draw[arrow] (sq2) -- (summe);
        \draw[arrow] (summe) -- (final);
    }{22}
}
}